\section{T-III and T-IV: Proposed Approaches}

Table~\ref{tab:status} separates completed evidence from proposed work. ``Not completed'' means
that no empirical result is claimed.

\begin{table}[t]
    \centering
    \caption{Assignment status at submission time.}
    \label{tab:status}
    \footnotesize
    \begin{tabular}{@{}p{0.13\linewidth}p{0.22\linewidth}p{0.53\linewidth}@{}}
        \toprule
        Stage & Status & Evidence / next gate \\
        \midrule
        T-I & Complete & Fixed 15k checkpoint; \timean\% $\pm$ \tistd{} pp over 3$\times$100 rollouts. \\
        T-II & Not completed & Tooling and protocol implemented; no valid discovery rollouts, two verified modes, targeted rates, or MP4 evidence. \\
        T-III & Not completed & Code scaffold only; no LLM reward or sampler was evaluated. \\
        T-IV & Not completed & No residual policy was trained or evaluated. \\
        \bottomrule
    \end{tabular}
\end{table}

\subsection{T-III: proposed LLM reward and episode generation}

Given a verified T-II rollout, the intended pipeline would provide the MP4, exact Push-T state
schema, official success definition, and executable simulator API to an LLM. The model would
return (1) a dense reward combining change in position error, wrapped orientation error, overlap
progress, action regularization, and terminal success, and (2) a sampler concentrated on the
measured failure-pose interval while retaining some nominal-distribution starts. Following
Eureka's code-generation and reward-reflection principle~\cite{ma2023eureka} and ROSETTA-style
staged grounding~\cite{rosetta}, the model would first describe the observed mechanism, then map
video claims to available state variables, and only then emit Python.

Generated code would be treated as an untrusted candidate. Static validation would reject unknown
state fields, side effects, non-finite rewards, out-of-domain reset bounds, and changes to the
success condition. Unit tests would evaluate each term on constructed states: moving toward the
goal should not reduce reward, orientation must use wrapped angles, terminal success should
dominate shaping, and inactivity must not be profitable. Short rollouts would then expose reward
hacking, scale imbalance, or an overly narrow sampler. The implemented scaffold preserves the
prompt, raw response, extracted code, validation report, and video hash, but no generated output
was experimentally evaluated.

\subsection{T-IV: proposed bounded residual PPO}

The intended T-IV method would freeze the RGB Diffusion Policy and learn only a bounded residual
with PPO:
\begin{equation}
    a_t = \operatorname{clip}\!\left(a_t^{\mathrm{DP}} + \alpha
    \tanh(a_t^{\mathrm{res}}), -1, 1\right).
\end{equation}
The residual would observe the same policy inputs plus state variables explicitly authorized for
the RL experiment; $\alpha$ would be conservative so the correction could not replace the base
controller wholesale. Training would use the audited T-III reward and a mixture of
failure-conditioned and nominal starts.

Evaluation would keep the T-I checkpoint, success threshold, horizons, and seeds fixed. Each
failure mode would compare base and residual policies on 100 targeted episodes at three seeds. A
second 3$\times$100 evaluation on the original Push-T distribution would measure catastrophic
forgetting. Necessary ablations include base only, residual without targeted sampling, targeted
sampling with sparse reward, and the full method. No PPO curve or improvement is reported because
this experiment was not run.
